\documentclass[12pt]{article}
\usepackage[T1]{fontenc}
\usepackage[utf8]{inputenc}
\usepackage{lmodern}
\usepackage{textcomp}
\usepackage{enumitem}
\usepackage{geometry}
\geometry{margin=1in}
\begin{document}
\begin{center}
Answer Key - Business Administration - Undergraduate Level Assessment 2 \
\small (Answers are ordered exactly as in the question paper.)
\end{center}

\section*{Multiple Choice}
\begin{enumerate}[label=\arabic*.]
\item \textbf{Answer:} D
\\ \textit{Options:} A) Contingency planning; B) Scenario planning; C) Financial planning; D) Building planning
\\ \textit{Note:} Building planning is not a recognized form of business continuity planning because business continuity planning focuses on strategies and procedures to ensure critical business operations can continue during and after a disruption, rather than on physical infrastructure development.
\item \textbf{Answer:} C
\\ \textit{Options:} A) High involvement buying situation.; B) New buying situation.; C) Routine buying.; D) Impulsive buying.
\\ \textit{Note:} Routine buying requires the least effort because it involves low-involvement decisions for frequently purchased items, where consumers rely on established habits and brand familiarity rather than extensive information search or evaluation.
\item \textbf{Answer:} C
\\ \textit{Options:} A) Search marketing.; B) Internet advertising.; C) Permission-based email marketing.; D) Social web marketing.
\\ \textit{Note:} Permission-based email marketing is correct because it specifically refers to the practice of sending marketing emails to recipients who have explicitly agreed to receive such communications, ensuring compliance with anti-spam laws and fostering a more engaged audience.
\item \textbf{Answer:} B
\\ \textit{Options:} A) Situation analysis; B) Strategy; C) Implementation; D) Evaluation involves assessing all aspects of the programme.
\\ \textit{Note:} The correct answer is "Strategy" because this stage of the communications plan focuses on determining the key messages and the rationale behind them, ensuring that communication efforts align with organizational goals and target audience needs.
\item \textbf{Answer:} D
\\ \textit{Options:} A) Matrix; B) Divisional; C) Multi-divisional; D) Functional
\\ \textit{Note:} Specialisation is a feature of the functional organisational structure because it allows for the grouping of employees based on their specific skills and expertise, leading to increased efficiency and productivity within distinct departments.
\end{enumerate}

\section*{True / False}
\begin{enumerate}[label=\arabic*.]
\setcounter{enumi}{5}
\item \textbf{Answer:} False
\\ \textit{Options:} A) True; B) False
\\ \textit{Note:} The ethical factor is not a competing responsibility but rather a foundational aspect that underpins corporate social responsibility, guiding companies in their commitment to ethical practices and stakeholder interests.
\item \textbf{Answer:} False
\\ \textit{Options:} A) True; B) False
\\ \textit{Note:} The statement is false because effective marketing communications prioritize the development of compelling messages that resonate with the target audience, ultimately influencing their perceptions and decisions rather than merely making products available.
\item \textbf{Answer:} False
\\ \textit{Options:} A) True; B) False
\\ \textit{Note:} Henry Ford is primarily recognized for his innovations in the automotive industry and the development of assembly line production, rather than for his philanthropic efforts in education and public health initiatives.
\item \textbf{Answer:} True
\\ \textit{Options:} A) True; B) False
\\ \textit{Note:} Public relations plays a crucial role in shaping and maintaining a company's reputation by addressing negative publicity and actively promoting positive narratives to enhance the corporate image.
\item \textbf{Answer:} False
\\ \textit{Options:} A) True; B) False
\\ \textit{Note:} Lateral communication refers to the exchange of information among individuals at the same organizational level, rather than vertical communication that occurs up and down the hierarchy.
\end{enumerate}

\section*{Long Form Response}
\begin{enumerate}[label=\arabic*.]
\setcounter{enumi}{10}
\item \textbf{Answer:} In business administration, various research methodologies are employed in evaluations, including quantitative, qualitative, and mixed-methods approaches. Quantitative research focuses on numerical data and statistical analysis, while qualitative research emphasizes understanding behaviors and motivations through interviews and observations. Mixed-methods combine both approaches to provide a comprehensive view of a research problem. However, behavior study is often not included in these methodologies as it tends to focus more on psychological aspects and individual actions, which may not align with the organizational and operational focus typically seen in business evaluations. This distinction highlights the importance of aligning research methodologies with specific business objectives and contexts.
\\ \textit{Note:} correct because it accurately describes the primary research methodologies in business administration-quantitative, qualitative, and mixed-methods-while also noting that behavior study is usually excluded due to its focus on individual behaviors rather than systematic data analysis relevant to business evaluations.
\item \textbf{Answer:} Improving a website's structure and content to enhance visibility in organic search engine results is fundamentally rooted in search engine optimization (SEO). Key principles include ensuring a clear site architecture that promotes easy navigation, as this helps search engines crawl and index content efficiently. Additionally, incorporating relevant keywords strategically throughout the website-within titles, headings, and body text- aligns the content with user search intents, thereby increasing the likelihood of higher rankings. Furthermore, regularly updating content and ensuring it is high-quality, informative, and engaging can lead to improved user retention and lower bounce rates, which positively influence search engine rankings. Overall, effective SEO combines structural integrity with keyword relevance to create a compelling online presence.
\\ \textit{Note:} correct because it emphasizes the importance of a clear site structure for effective crawling and indexing by search engines, as well as the strategic use of relevant keywords to align website content with user search intents, both of which are essential components of effective SEO practices.
\end{enumerate}

\vfill
\noindent\textit{End of Answer Key}
\end{document}